\section{Optimism}
\label{sec:optimism}

Optimism es una red de capa 2 de Ethereum del tipo \textit{rollup} optimista, construida sobre
el OP Stack, un conjunto de componentes de código abierto que permite desplegar cadenas con la
misma arquitectura~\cite{opstackoverview}. Su red de prueba, Optimism Sepolia, se apoya sobre la
red de prueba Ethereum Sepolia y actúa en este trabajo como cadena de destino del puente, es decir,
como la cadena en la que se despliegan el contrato verificador de las pruebas de conocimiento
cero y el contrato que acuña los activos. La presente sección describe la arquitectura de
Optimism, el ciclo de confirmación de sus bloques, su modelo de tarifas, el mecanismo de pruebas
de fraude que sostiene su seguridad y el puente nativo que lo comunica con Ethereum, con el fin
de precisar las propiedades de la cadena de destino que condicionan el diseño del puente.

\subsection{Arquitectura del OP Stack}
\label{sec:op_arquitectura}

De acuerdo con la clasificación presentada en la Sección~\ref{sec:aztec_rollups}, Optimism
ejecuta las transacciones fuera de Ethereum, publica en Ethereum los datos necesarios para
reconstruir su estado y presume la validez de los estados resultantes mientras nadie demuestre lo
contrario. Su especificación la describe como una cadena \textit{derivada}, puesto que su
historia canónica queda definida exclusivamente por los datos publicados en la capa
1~\cite{opstackoverview}. Cualquier nodo que lea dichos datos desde Ethereum y aplique las reglas
de derivación obtiene la misma cadena que construyó el operador de la red.

Al igual que un nodo de Ethereum posterior a \textit{The Merge} (Sección~\ref{sec:eth_capas}), un
nodo de Optimism se compone de un cliente de consenso, denominado nodo de \textit{rollup}, y de un
cliente de ejecución derivado de un cliente de Ethereum, ambos comunicados mediante la misma
interfaz que utilizan los clientes de la capa 1. Sobre esta base, el OP Stack distingue cuatro
funciones principales~\cite{opstackoverview}:

\begin{itemize}
	\item El \textbf{secuenciador} recibe las transacciones de los usuarios, las ordena y produce
	los bloques de la capa 2. En Optimism Sepolia se produce un bloque cada dos
	segundos~\cite{superchainopsepolia}.
	\item El \textbf{publicador de lotes} (\textit{batcher}) comprime los bloques producidos y
	publica sus datos en Ethereum, habitualmente en forma de \textit{blobs}, un tipo de dato
	introducido por la propuesta EIP-4844 que ofrece un almacenamiento temporal más económico que
	el \textit{calldata}~\cite{eip4844}.
	\item El \textbf{proponente} (\textit{proposer}) publica periódicamente en Ethereum
	\textit{raíces de salida} (\textit{output roots}), que son compromisos sobre el estado
	completo de la capa 2 en un bloque determinado.
	\item El \textbf{impugnador} (\textit{challenger}) vigila las raíces de salida publicadas y
	las disputa cuando no coinciden con el estado que resulta de derivar la cadena.
\end{itemize}

En la capa 1, estas funciones se apoyan en un conjunto de contratos inteligentes. Entre ellos
destaca el contrato \texttt{OptimismPortal}, que recibe los depósitos y procesa los retiros.
También destaca el contrato \texttt{DisputeGameFactory}, que registra las raíces de salida y crea
los juegos de disputa asociados a ellas~\cite{opspecsoverview}. La Figura~\ref{fig:op_arquitectura} representa la
relación entre los componentes de ambas capas.

\begin{figure}[htbp]
	\centering
	\begin{tikzpicture}[
		caja/.style={draw, rounded corners=2pt, align=center, minimum width=3.6cm,
			minimum height=1.1cm, font=\small},
		capa/.style={draw, dashed, rounded corners=4pt, inner sep=0.3cm},
		flecha/.style={-{Stealth[length=2.2mm]}, thick},
		etiqueta/.style={font=\footnotesize\itshape, align=center}
		]
		% Capa 2
		\node[caja] (seq) at (0,0) {Secuenciador\\ (bloques cada 2~s)};
		\node[caja] (bat) at (4.8,0) {Publicador\\ de lotes};
		\node[caja] (prop) at (9.6,0) {Proponente};
		% Capa 1
		\node[caja] (portal) at (0,-3.6) {\texttt{OptimismPortal}\\ (depósitos y retiros)};
		\node[caja] (blobs) at (4.8,-3.6) {Datos de lotes\\ (\textit{blobs})};
		\node[caja] (dgf) at (9.6,-3.6) {\texttt{DisputeGameFactory}\\ (raíces de salida)};

		\node[capa, fit=(seq)(bat)(prop), label={[font=\small\bfseries]above:Optimism (capa 2)}] {};
		\node[capa, fit=(portal)(blobs)(dgf), label={[font=\small\bfseries]below:Ethereum (capa 1)}] {};

		\draw[flecha] (seq) -- node[etiqueta, above] {bloques} (bat);
		\draw[flecha] (bat) -- node[etiqueta, above] {estado} (prop);
		\draw[flecha] (portal) -- node[etiqueta, right] {depósitos} (seq);
		\draw[flecha] (bat) -- node[etiqueta, right] {lotes\\ comprimidos} (blobs);
		\draw[flecha] (prop) -- node[etiqueta, right] {raíces\\ de salida} (dgf);
		\draw[flecha] (dgf) -- node[etiqueta, above] {raíz aceptada} (portal);
	\end{tikzpicture}
	\caption{Componentes del OP Stack y su relación con los contratos desplegados en Ethereum.}
	\label{fig:op_arquitectura}
\end{figure}

El secuenciador no puede excluir de forma indefinida los depósitos realizados en la capa 1. Las
transacciones enviadas a \texttt{OptimismPortal} deben incorporarse a la capa 2 dentro de una
ventana de secuenciación, que en Optimism Sepolia es de 3\,600 bloques de
Ethereum~\cite{superchainopsepolia}, y en caso contrario las reglas de derivación las incluyen de
todas formas~\cite{opspecsoverview}. Esta propiedad otorga a los usuarios resistencia a la
censura aun cuando el secuenciador sea un operador único.

\subsection{Confirmación de bloques y finalidad}
\label{sec:op_finalidad}

Dado que los bloques de Optimism se producen antes de que sus datos lleguen a Ethereum, un mismo
bloque atraviesa distintos estados de confirmación a lo largo del tiempo. La
Tabla~\ref{tab:op_estados} resume los tres estados que distingue el OP
Stack~\cite{opstackoverview}.

\begin{table}[htbp]
	\centering
	\caption{Estados de confirmación de un bloque de Optimism.}
	\label{tab:op_estados}
	\small
	\begin{tabular}{>{\raggedright\arraybackslash}p{2.6cm} >{\raggedright\arraybackslash}p{6.2cm} >{\raggedright\arraybackslash}p{5.4cm}}
		\toprule
		\textbf{Estado} & \textbf{Condición} & \textbf{Garantía} \\
		\midrule
		No seguro (\textit{unsafe}) & El secuenciador produjo el bloque y lo difundió por la red
		\textit{peer-to-peer}. & Depende de la honestidad del secuenciador. \\
		\addlinespace
		Seguro (\textit{safe}) & Los datos del bloque fueron publicados en Ethereum y el bloque
		puede derivarse a partir de ellos. & Depende de que el bloque de Ethereum que contiene los
		datos no sea revertido. \\
		\addlinespace
		Finalizado (\textit{finalized}) & El bloque de Ethereum que contiene los datos alcanzó la
		finalidad. & Equivalente a la finalidad de Ethereum
		(Sección~\ref{sec:eth_consenso}). \\
		\bottomrule
	\end{tabular}
\end{table}

Conviene distinguir este concepto de finalidad del que se aplica a los retiros hacia Ethereum.
Un bloque finalizado de Optimism no puede ser revertido, puesto que sus datos forman parte de la
historia finalizada de Ethereum y su contenido queda determinado por las reglas de derivación.
Sin embargo, Ethereum no ejecuta esas reglas por sí mismo, de modo que solo acepta una afirmación
sobre el estado de Optimism, como la existencia de un retiro, una vez que la raíz de salida
correspondiente ha superado el período de impugnación descrito en la
Sección~\ref{sec:op_fraude}.

\subsection{Compatibilidad con la EVM y modelo de tarifas}
\label{sec:op_tarifas}

Optimism busca la \textbf{equivalencia con la EVM}, lo que significa que su cliente de ejecución
reproduce el comportamiento de la máquina virtual de Ethereum con un número reducido de
modificaciones, documentadas en su especificación~\cite{opspecsexec}. Los contratos escritos para
Ethereum pueden desplegarse en Optimism sin cambios, y los contratos precompilados descritos en la
Sección~\ref{sec:eth_gas}, incluidos los de la curva BN254, se encuentran disponibles con los
mismos costos de gas. La actualización Isthmus incorporó, además, los precompilados de BLS12-381
definidos en la propuesta EIP-2537~\cite{opisthmus, eip2537}. Esta propiedad permite que el
contrato verificador generado por las herramientas de Noir y de SP1 para Ethereum se despliegue
en Optimism Sepolia sin modificaciones.

El modelo de tarifas, en cambio, presenta una diferencia relevante. Las transacciones de Optimism
pagan la tarifa base y la propina del mecanismo EIP-1559 por su ejecución en la capa 2, y pagan
además una \textbf{tarifa de datos de capa 1} (\textit{L1 data fee}), que cubre el costo de
publicar sus datos en Ethereum~\cite{opspecsexec}. Esta segunda tarifa depende del tamaño de la
transacción una vez comprimida y de los precios vigentes de \textit{blobs} y de gas en Ethereum,
que la capa 2 conoce a través del contrato predesplegado \texttt{L1Block}. En consecuencia, el
costo total de verificar una prueba en Optimism no depende solo del gas de ejecución, sino
también del tamaño de la prueba y de las entradas públicas que se transmiten en la transacción.
Por esta razón, en este trabajo el gas de ejecución consumido por la verificación se reporta por
separado del costo monetario, que varía con las condiciones de ambas capas.

\subsection{Pruebas de fraude}
\label{sec:op_fraude}

La seguridad de las raíces de salida de Optimism descansa en un sistema de pruebas de fraude
habilitado sin permisos en la red principal desde junio de 2024~\cite{optimismfaultproofs}.
Cuando un participante publica una raíz de salida, cualquier otro puede impugnarla depositando
una garantía, lo que da inicio a un \textbf{juego de disputa}. En él, ambas partes acotan de
forma iterativa, mediante bisección, el primer paso de la ejecución sobre el cual discrepan, hasta
reducir el desacuerdo a una única instrucción. Esa instrucción se ejecuta entonces en un contrato
de Ethereum que implementa la máquina virtual Cannon, la cual reproduce el programa de derivación
y ejecución de Optimism compilado para una arquitectura MIPS~\cite{opcannon}. El resultado de esta
ejecución determina qué parte tenía razón, y la parte equivocada pierde su garantía.

Si ninguna impugnación prospera durante el período previsto, de aproximadamente siete días, la
raíz de salida se considera válida y puede utilizarse para procesar retiros~\cite{optimismfaultproofs}.
El modelo de seguridad resultante se apoya en el supuesto de que existe al menos un participante
honesto que vigila la cadena y que puede publicar una impugnación en Ethereum dentro de ese plazo.
Este supuesto es más débil que el de confiar en un operador concreto, pero difiere de la
garantía que ofrece una prueba de validez, en la que la corrección del estado se verifica de forma
inmediata y no depende de la vigilancia de terceros.

\subsection{Puente nativo con Ethereum}
\label{sec:op_puente_nativo}

Optimism dispone de un puente nativo entre ambas capas, compuesto por los contratos
\texttt{L1CrossDomainMessenger} y \texttt{L1StandardBridge} en Ethereum y sus equivalentes
predesplegados en la capa 2~\cite{opspecsoverview}. En un \textbf{depósito}, el contrato de la
capa 1 bloquea los activos y emite un mensaje a través de \texttt{OptimismPortal}, que la capa 2
incorpora como una transacción especial capaz de acuñar los activos equivalentes. Como el
contenido de la capa 2 se deriva de los datos de Ethereum, el depósito queda acreditado en pocos
minutos y sin intervención de terceros.

Un \textbf{retiro} sigue el camino inverso y requiere dos pasos en
Ethereum~\cite{opspecswithdrawals}. Primero, el usuario \textit{prueba} el retiro, demostrando
mediante una prueba de Merkle que el mensaje fue registrado en el contrato
\texttt{L2ToL1MessagePasser} y que este registro forma parte de una raíz de salida publicada.
Después, una vez transcurrido el período de impugnación y resuelto a favor el juego de disputa
correspondiente, el usuario \textit{finaliza} el retiro y el contrato libera los activos
bloqueados. La asimetría entre ambos sentidos refleja la naturaleza del \textit{rollup}
optimista, puesto que la capa 2 conoce directamente el estado de Ethereum, mientras que Ethereum
solo conoce el estado de Optimism a través de afirmaciones que deben superar el período de
impugnación.

La capa 2 conoce el estado de Ethereum porque cada bloque de Optimism comienza con una
transacción del sistema que registra en el contrato \texttt{L1Block} el número, la marca de
tiempo y el hash del bloque de Ethereum que le sirve de origen~\cite{opspecspredeploys}. El
puente desarrollado en este trabajo no se apoya en este mecanismo, que es propio de la relación
entre un \textit{rollup} y su capa 1, sino en la verificación de la firma del comité de
sincronización de Ethereum descrita en la Sección~\ref{sec:eth_sync_committee}. De este modo, la
solución resulta aplicable a cualquier cadena compatible con la EVM, con independencia de que
mantenga o no una relación nativa con Ethereum, y Optimism Sepolia cumple el papel de entorno de
validación representativo de una cadena de destino con costos de ejecución reducidos.
